
%\needspace{4\baselineskip}
\color{uniblue}{\subsection{Защита от непереключения фаз (ЗНФ)}}
\color{black}

\begin{enumerate}[label=\arabic{section}.\arabic{subsection}.\arabic{enumi}, labelsep=4pt, leftmargin=0pt, itemindent=57pt, itemsep=0pt, parsep=0pt] % если что можно крутить itemsep 0 или 5

\item
Защита от непереключения фаз (ЗНФ) предназначена для обнаружения расхождения полюсов выключателя с пофазным приводом при его включении или отключении. 

\item
Ввод в работу ЗНФ выключателя В1 осуществляется программным ключом \textbf{«ЗНФ В1»}. Факт расхождения полюсов может устанавливаться одним из следующих способов:

\begin{enumerate}[label=\arabic*)]
\item на основании внешнего сигнала от сборки блок-контактов выключателя;
\item функцией фиксации положения выключателя (принцип работы описан в пункте 1.20 настоящего руководства по эксплуатации), которая пофазно контролирует состояние блок-контактов выключателя и в случае расхождения полюсов формирует сигнал \textit{«Внутр. пуск. ЗНФ В1»}.
\end{enumerate}

\item
В зависимости от используемого способа соответствующий сигнал должен назначаться на конфигурируемый вход ФСУ \textit{«Пуск ЗНФ В1»}.

\item
Действие на отключение выключателя В1 осуществляется с выдержкой времени \textbf{«Тср ЗНФ В1»}. Задаваемая выдержка времени предназначена для отстройки от разновременности переключения блок-контактов выключателя.

\item
Предусмотрена возможность блокировки ЗНФ В1 от внешнего сигнала \textit{«Блок. ЗНФ В1»} на время, которое задается уставкой \textbf{«Тдеблок ЗНФ В1»}.

\item
ЗНФ В1 может быть оперативно выведена из работы от внешнего сигнала \textit{«Вывод ЗНФ В1»}, назначенного на соответствующий дискретный вход или функциональную клавишу.

\item
В устройстве также предусмотрена возможность реализации функции ЗНФ выключателя В2 -- вводится программным ключом \textbf{«ЗНФ В2»}. Действие на отключение выключателя В2 осуществляется с выдержкой времени \textbf{«Тср ЗНФ В2»}. Функция ЗНФ В2 оперативно выводится от внешнего сигнала \textit{«Вывод ЗНФ В2»}.

\item
Функциональная схема логики ЗНФ В1 приведена на рисунке \ref{pict_bk:ZNF}, функциональная схема логики ЗНФ В2 выполняется аналогично.

\begin{figure}[H]
\centering
\includegraphics[width=1\textwidth,height=1\textheight,keepaspectratio]{\fbpath/761.3441 ЗНФ/img/ZNF.pdf}
\captionof{figure}{ Функциональная схема ЗНФ В1}\label{pict_bk:ZNF}
\end{figure}

\item
Уставки ЗНФ представлены в таблице \ref{desc_znf:tbl1}.

\small
\begin{longtable}{|>{\centering\arraybackslash}m{5.3cm}|>{\centering\arraybackslash}m{3.3cm}|>{\centering\arraybackslash}m{4.2cm}|>{\centering\arraybackslash}m{1.8cm}|>{\centering\arraybackslash}m{1cm}|}
\caption{Параметры для настройки ЗНФ\hfill\vspace{-0.5\baselineskip}}\label{desc_znf:tbl1}\\ 
\hline
\rowcolor{gray!30}
Параметр & Условное обозначение на схеме & Значение/ Диапазон & Единица измерения & Шаг \\ 
\hline
\endfirsthead
\caption*{\hspace{3pt}\emph{Продолжение таблицы \ref{desc_znf:tbl1}\hfill\vspace{-0.5\baselineskip}}} \\ % сделано по ГОСТ 2.105 п.6.8.7
\hline
\rowcolor{gray!30}
Параметр & Условное обозначение на схеме & Значение/ Диапазон & Единица измерения & Шаг \\ 
%\hline
\endhead

%===m>ЗНФ В1|ЗНФ
\raggedright {'Description': ''} & \centering Тср ЗНФ В & \centering 0 ... 100000 & \centering  & \centering \arraybackslash 1 \\
\hline
\raggedright {'Description': ''} & \centering Тдеблок ЗНФ В & \centering 0 ... 100000 & \centering  & \centering \arraybackslash 1 \\
\hline
\raggedright {'Description': ''} & \centering ЗНФ В1 & \centering Выведено \\ Введено & \centering --- & \centering \arraybackslash --- \\
\hline
%===m
\end{longtable}
\normalsize

\end{enumerate}